% sec2_3.tex — Section 2.3: Large-Sample Distribution of the OLS Estimator

The importance in econometrics of the OLS procedure, originally developed for the
classical regression model of Chapter~1, lies in the fact that it has good asymptotic
properties for a class of models, different from the classical model, that are useful
in economics.  Of those models, the model presented in this section has probably the
widest range of economic applications.  No specific distributional assumption (such
as the normality of the error term) is required to derive the asymptotic distribution
of the OLS estimator.  The requirement in finite-sample theory that the regressors be
strictly exogenous or ``fixed'' is replaced by a much weaker requirement that they be
``predetermined.''  (For the sake of completeness, the appendix develops a parallel
asymptotic theory for a model with ``fixed'' regressors.)

%% ─── The Model ─────────────────────────────────────────────────────────────
\subsection*{The Model}

We use the term the \textbf{data generating process} (\textbf{DGP}) for the stochastic
process that generated the finite sample $(\mathbf{y}, \mathbf{X})$.  Therefore, if we
specify the DGP, the joint distribution of the finite sample $(\mathbf{y}, \mathbf{X})$
can be determined.  In finite-sample theory, where the sample size is fixed and
finite, we defined a model as a set of the joint distributions of
$(\mathbf{y}, \mathbf{X})$.  In large-sample theory, a model is stated as a set of
DGPs.  The model we study is the set of DGPs satisfying the following set of
assumptions.

\begin{assumption}[linearity]
\label{ass:2.1}
\begin{equation*}
  y_i = \mathbf{x}_i'\bm{\beta} + \varepsilon_i
  \quad (i = 1, 2, \ldots, n)
\end{equation*}
\emph{where $\mathbf{x}_i$ is a $K$-dimensional vector of explanatory variables
(regressors), $\bm{\beta}$ is a $K$-dimensional coefficient vector, and
$\varepsilon_i$ is the unobservable error term.}
\end{assumption}

\begin{assumption}[ergodic stationarity]
\label{ass:2.2}
\emph{The $(K+1)$-dimensional vector stochastic process $\{y_i, \mathbf{x}_i\}$ is
jointly stationary and ergodic.}
\end{assumption}

\begin{assumption}[predetermined regressors]
\label{ass:2.3}
\emph{All the regressors are \textbf{predetermined} in the sense that they are
orthogonal to the contemporaneous error term:
$\E(x_{ik}\varepsilon_i) = 0$ for all $i$ and $k$ $(= 1, 2, \ldots, K)$.\footnote{Our
definition of the term \emph{predetermined} is not universal.  Some authors say the
regressors are predetermined if
$\E(\mathbf{x}_{i-j} \cdot \varepsilon_i) = \mathbf{0}$ for all $j \geq 0$, not just
for $j = 0$.  That is, the error term is orthogonal not only to the contemporaneous
but also the past regressors.  In Koopmans and Hood (1953), the regressors are said
to be predetermined if $\varepsilon_i$ is independent of $\mathbf{x}_{i-j}$ for all
$j \geq 0$.  Our definition is the same as Hamilton's (1994).}
This can be written as}
\[
  \E[\mathbf{x}_i \cdot (y_i - \mathbf{x}_i'\bm{\beta})] = \mathbf{0}
  \quad \textit{or equivalently} \quad
  \E(\mathbf{g}_i) = \mathbf{0}
\]
\emph{where $\mathbf{g}_i \equiv \mathbf{x}_i \cdot \varepsilon_i$.}
\end{assumption}

\begin{assumption}[rank condition]
\label{ass:2.4}
\emph{The $K \times K$ matrix $\E(\mathbf{x}_i\mathbf{x}_i')$ is nonsingular (and
hence finite).  We denote this matrix by $\bm{\Sigma}_{\mathbf{xx}}$.}
\end{assumption}

\begin{assumption}[$\mathbf{g}_i$ is a martingale difference sequence with finite
second moments]
\label{ass:2.5}
\emph{$\{\mathbf{g}_i\}$ is a martingale difference sequence (so a fortiori
$\E(\mathbf{g}_i) = \mathbf{0}$).  The $K \times K$ matrix of cross moments,
$\E(\mathbf{g}_i\mathbf{g}_i')$, is nonsingular.  We use $\mathbf{S}$ for
$\avar(\bar{\mathbf{g}})$ (the variance of the asymptotic distribution of
$\sqrt{n}\,\bar{\mathbf{g}}$, where
$\bar{\mathbf{g}} \equiv \frac{1}{n}\sum_i \mathbf{g}_i$).  By Assumption~\ref{ass:2.2}
and the ergodic stationary Martingale Differences CLT,
$\mathbf{S} = \E(\mathbf{g}_i\mathbf{g}_i')$.}
\end{assumption}

The first assumption is just reproducing Assumption~1.1.  The rest of the assumptions
require some lengthy comments.

\begin{itemize}
\item (Ergodic stationarity)\quad A trivial but important special case of ergodic
  stationarity is that $\{y_i, \mathbf{x}_i\}$ is i.i.d., that is, the sample is a
  random sample.\footnote{Actually, once the independence assumption is made, the same
  large-sample results can be proved for the more general case where
  $\{y_i, \mathbf{x}_i\}$ is independently but not identically distributed (i.n.i.d.),
  provided that some conditions on higher moments of the joint distribution of
  $(\varepsilon_i, \mathbf{x}_i)$ are satisfied.  We will not entertain this
  generalization because the i.i.d.\ assumption is satisfied in most microdata sets.}
  Most existing microdata on households are random samples, with observations
  randomly drawn from a population of a nation's households.  Thus, we are in no
  way ruling out models that use cross-section data.

\item (The model accommodates conditional heteroskedasticity)\quad
  If $\{y_i, \mathbf{x}_i\}$ is stationary, then the error term
  $\varepsilon_i = y_i - \mathbf{x}_i'\bm{\beta}$ is also stationary.  Thus,
  Assumption~\ref{ass:2.2} implies that the unconditional second moment
  $\E(\varepsilon_i^2)$---if it exists and is finite---is constant across $i$.  That
  is, the error term is unconditionally homoskedastic.  Yet the error can be
  conditionally heteroskedastic in that the conditional second moment,
  $\E(\varepsilon_i^2 \mid \mathbf{x}_i)$, can depend on $\mathbf{x}_i$.  An example
  in which the error is homoskedastic unconditionally but not conditionally is
  included in Section~2.6, where the consequence of superimposing conditional
  homoskedasticity (that $\E(\varepsilon_i^2 \mid \mathbf{x}_i) = \sigma^2$) on the
  model will be explored.

\item ($\E(\mathbf{x}_i \cdot \varepsilon_i) = \mathbf{0}$ vs.\
  $\E(\varepsilon_i \mid \mathbf{x}_i) = 0$)\quad Sometimes, instead of the
  orthogonality condition $\E(\mathbf{x}_i \cdot \varepsilon_i) = \mathbf{0}$, it is
  assumed that the error is unrelated in the sense that
  $\E(\varepsilon_i \mid \mathbf{x}_i) = 0$.  This is stronger than the orthogonality
  condition because it implies that, for any (measurable) function $f$ of
  $\mathbf{x}_i$, $f(\mathbf{x}_i)$ is orthogonal to $\varepsilon_i$:
  \[
    \E[f(\mathbf{x}_i)\varepsilon_i]
    = \E[\E(f(\mathbf{x}_i)\varepsilon_i \mid \mathbf{x}_i)]
    = \E[f(\mathbf{x}_i)\,\E(\varepsilon_i \mid \mathbf{x}_i)] = 0.
  \]
  In rational expectations models, this stronger condition is satisfied, but for the
  purpose of developing asymptotic theory, we need only the weaker assumption of the
  orthogonality condition.

\item (Predetermined vs.\ strictly exogenous regressors)\quad The regressors are
  \emph{not} required to be strictly exogenous.  As we noted in Section~1.2, the
  exogeneity assumption (Assumption~1.2) implies that, for any regressor $k$,
  $\E(x_{jk}\varepsilon_i) = 0$ for all $i$ and $j$, not just for $i = j$, which
  rules out the possibility that the current error term, $\varepsilon_i$, is
  correlated with future regressors, $\mathbf{x}_{i+j}$ for $j \geq 1$.
  Assumption~\ref{ass:2.3}, restricting only the \emph{contemporaneous} relationship
  between the error term and the regressors, does not rule out that possibility.
  For example, the AR(1) process, which does not satisfy the exogeneity assumption of
  the classical regression model, can be accommodated in the model of this chapter.
  This weaker assumption of predetermined regressors will be further relaxed in the
  next chapter.

\item (Rank condition as no multicollinearity in the limit)\quad
  Since $\E(\mathbf{x}_i\mathbf{x}_i')$ is finite by Assumption~\ref{ass:2.4},
  $\lim_{n \to \infty} \mathbf{S}_{\mathbf{xx}} = \bm{\Sigma}_{\mathbf{xx}}$
  (where $\mathbf{S}_{\mathbf{xx}} \equiv \frac{1}{n}\sum_{i=1}^{n}
  \mathbf{x}_i\mathbf{x}_i'$) with probability one by the Ergodic Theorem.  So, for
  $n$ sufficiently large, the sample cross moment of the regressors
  $\mathbf{S}_{\mathbf{xx}}$, which can be written as
  $\frac{1}{n}\mathbf{X}'\mathbf{X}$, is nonsingular by Assumptions~\ref{ass:2.2} and
  \ref{ass:2.4}.  Since $\frac{1}{n}\mathbf{X}'\mathbf{X}$ is nonsingular if and only
  if $\rank(\mathbf{X}) = K$, Assumption~1.3 (no multicollinearity) is satisfied with
  probability one for sufficiently large $n$.  In the OLS formula
  $\mathbf{b} = \mathbf{S}_{\mathbf{xx}}^{-1}\mathbf{s}_{\mathbf{xy}}$
  (where $\mathbf{s}_{\mathbf{xy}} \equiv \frac{1}{n}\sum_{i=1}^{n}
  \mathbf{x}_i \cdot y_i$), $\mathbf{S}_{\mathbf{xx}}$ needs to be inverted.  If
  $\mathbf{S}_{\mathbf{xx}}$ is singular in a finite sample (so it cannot be
  inverted), we just assign an arbitrary value to $\mathbf{b}$ so that the OLS
  estimator is well-defined for any sample.

\item (A sufficient condition for $\{\mathbf{g}_i\}$ to be an m.d.s.)\quad
  Since an m.d.s.\ (martingale difference sequence) is zero-mean by definition,
  Assumption~\ref{ass:2.5} is stronger than Assumption~\ref{ass:2.3}.  We will need
  Assumption~\ref{ass:2.5} to prove the asymptotic normality of the OLS estimator.
  The assumption, about the product of the regressors and the error term, may be hard
  to interpret.  A sufficient condition that is easier to interpret is
  \begin{equation}
    \E(\varepsilon_i \mid \varepsilon_{i-1}, \varepsilon_{i-2}, \ldots,
    \varepsilon_1, \mathbf{x}_i, \mathbf{x}_{i-1}, \ldots, \mathbf{x}_1) = 0.
    \label{eq:2.3.1}
  \end{equation}
  Note that the \emph{current} as well as lagged regressors is included in the
  information set.  This condition implies that the error term is serially uncorrelated
  and also is uncorrelated with the current and past regressors (the proof is much
  like the proof in the previous section that an m.d.s.\ is serially uncorrelated).
  That \eqref{eq:2.3.1} is sufficient for $\{\mathbf{g}_i\}$ to be an m.d.s.\ can be
  seen as follows.  We have
  \begin{align}
    \E(\mathbf{g}_i \mid \mathbf{g}_{i-1}, \ldots, \mathbf{g}_1)
    &= \E[\E(\mathbf{g}_i \mid \varepsilon_{i-1}, \varepsilon_{i-2}, \ldots,
       \varepsilon_1, \mathbf{x}_i, \mathbf{x}_{i-1}, \ldots, \mathbf{x}_1)
       \mid \mathbf{g}_{i-1}, \ldots, \mathbf{g}_1].
       \notag
  \end{align}
  This holds by the Law of Iterated Expectations because there is more information
  in the ``inside'' information set
  $(\varepsilon_{i-1}, \varepsilon_{i-2}, \ldots, \varepsilon_1, \mathbf{x}_i,
  \mathbf{x}_{i-1}, \ldots, \mathbf{x}_1)$ than in the ``outside'' information set
  $(\mathbf{g}_{i-1}, \ldots, \mathbf{g}_1)$.  Therefore,
  \begin{align}
    \E(\mathbf{g}_i \mid \mathbf{g}_{i-1}, \ldots, \mathbf{g}_1)
    &= \E[\mathbf{x}_i\,\E(\varepsilon_i \mid \varepsilon_{i-1},
       \varepsilon_{i-2}, \ldots, \varepsilon_1, \mathbf{x}_i,
       \mathbf{x}_{i-1}, \ldots, \mathbf{x}_1)
       \mid \mathbf{g}_{i-1}, \ldots, \mathbf{g}_1] \notag \\
    &\quad \text{(by the linearity of conditional expectations)} \notag \\
    &= \mathbf{0} \quad (\text{by } \eqref{eq:2.3.1}).
    \label{eq:2.3.2}
  \end{align}

\item (When the regressors include a constant)\quad In virtually all applications,
  the regressors include a constant so that $x_{i1} = 1$ for all $i$, then
  Assumption~\ref{ass:2.3} of predetermined regressors can be stated in more familiar
  terms: the mean of the error term is zero (which is implied by
  $\E(x_{ik}\varepsilon_i) = 0$ for $k = 1$), and the contemporaneous correlation
  between the error term and the regressors is zero (which is implied by
  $\E(x_{ik}\varepsilon_i)$ for $k \neq 1$ and $\E(\varepsilon_i) = 0$).  Also,
  since the first element of the $K$-dimensional vector
  $\mathbf{g}_i$ $(\equiv \mathbf{x}_i \cdot \varepsilon_i)$ is $\varepsilon_i$,
  Assumption~\ref{ass:2.5} implies
  \[
    \E(\varepsilon_i \mid \mathbf{g}_{i-1}, \mathbf{g}_{i-2}, \ldots,
    \mathbf{g}_1) = 0.
  \]
  Then, by the Law of Iterated Expectations, $\{\varepsilon_i\}$ is a scalar m.d.s.\ and
  hence is serially uncorrelated.
  \begin{equation}
    \E(\varepsilon_i \mid \varepsilon_{i-1}, \varepsilon_{i-2}, \ldots,
    \varepsilon_1) = 0.
    \label{eq:2.3.3}
  \end{equation}
  Therefore, Assumption~\ref{ass:2.5} implies that the error term itself is an m.d.s.\
  and hence is serially uncorrelated.

\item ($\mathbf{S}$ is a matrix of fourth moments)\quad Since
  $\mathbf{g}_i \equiv \mathbf{x}_i \cdot \varepsilon_i$, the $\mathbf{S}$ in
  Assumption~\ref{ass:2.5} can be written as
  $\E(\varepsilon_i^2\mathbf{x}_i\mathbf{x}_i')$.  Its $(k,j)$ element is
  $\E(\varepsilon_i^2 x_{ik} x_{ij})$.  So $\mathbf{S}$ is a matrix of fourth moments
  (the expectation of products of four different variables).  Consistent estimation of
  $\mathbf{S}$ will require an additional assumption to be specified in Section~2.5.

\item ($\mathbf{S}$ will take a different expression without Assumption~\ref{ass:2.5})
  \quad Thanks to the assumption that $\{\mathbf{g}_i\}$ is an m.d.s.,
  $\mathbf{S}$ $(\equiv \avar(\bar{\mathbf{g}}))$ is equal to
  $\E(\mathbf{g}_i\mathbf{g}_i')$.  Without the assumption, as we will see in
  Chapter~6, the expression for $\mathbf{S}$ is more complicated and involves
  autocovariances of $\mathbf{g}_i$.
\end{itemize}

%% ─── Asymptotic Distribution of the OLS Estimator ─────────────────────────
\subsection*{Asymptotic Distribution of the OLS Estimator}

We now prove that the OLS estimator is consistent and asymptotically normal.  It
should be kept in mind throughout the rest of this chapter that the OLS estimator
$\mathbf{b}$ depends on the sample size $n$ (although the dependence is not made
explicit by our choice of not to subscript $\mathbf{b}$ by $n$) and that $K$, the
number of regressors, is held fixed when we track the sequence of OLS estimators
indexed by $n$.  For the time being, we presume that there is available some
consistent estimator, denoted $\widehat{\mathbf{S}}$, of
$\mathbf{S}$ $(\equiv \avar(\bar{\mathbf{g}}) = \E(\mathbf{g}_i\mathbf{g}_i')
= \E(\varepsilon_i^2\mathbf{x}_i\mathbf{x}_i'))$.  The issue of estimating
$\mathbf{S}$ consistently will be taken up later.

\begin{proposition}[asymptotic distribution of the OLS Estimator]
\label{prop:2.1}
\begin{enumerate}[label=(\alph*)]
  \item \emph{(Consistency of\/ $\mathbf{b}$ for $\bm{\beta}$)\quad Under
    Assumptions~\ref{ass:2.1}--\ref{ass:2.4},
    $\plim_{n \to \infty} \mathbf{b} = \bm{\beta}$.  (So Assumption~\ref{ass:2.5}
    is not needed for consistency.)}
  \item \emph{(Asymptotic Normality of\/ $\mathbf{b}$)\quad If
    Assumption~\ref{ass:2.3} is strengthened as Assumption~\ref{ass:2.5}, then}
    \[
      \sqrt{n}(\mathbf{b} - \bm{\beta}) \dto
      N(\mathbf{0},\, \avar(\mathbf{b}))
      \quad \text{as } n \to \infty,
    \]
    \emph{where}
    \begin{equation}
      \avar(\mathbf{b})
      = \bm{\Sigma}_{\mathbf{xx}}^{-1}\,\mathbf{S}\,
        \bm{\Sigma}_{\mathbf{xx}}^{-1}.
      \label{eq:2.3.4}
    \end{equation}
    \emph{(Recall: $\bm{\Sigma}_{\mathbf{xx}} \equiv \E(\mathbf{x}_i\mathbf{x}_i')$,
    $\mathbf{S} = \E(\mathbf{g}_i\mathbf{g}_i')$,
    $\mathbf{g}_i \equiv \mathbf{x}_i \cdot \varepsilon_i$.)}
  \item \emph{(Consistent Estimate of\/ $\avar(\mathbf{b})$)\quad Suppose there is
    available a consistent estimator, $\widehat{\mathbf{S}}$, of\/
    $\mathbf{S}$ $(K \times K)$.  Then, under Assumption~\ref{ass:2.2},
    $\avar(\mathbf{b})$ is consistently estimated by}
    \begin{equation}
      \widehat{\avar(\mathbf{b})}
      = \mathbf{S}_{\mathbf{xx}}^{-1}\,\widehat{\mathbf{S}}\,
        \mathbf{S}_{\mathbf{xx}}^{-1},
      \label{eq:2.3.5}
    \end{equation}
    \emph{where $\mathbf{S}_{\mathbf{xx}}$ is the sample mean of\/
    $\mathbf{x}_i\mathbf{x}_i'$:}
    \begin{equation}
      \underset{(K \times K)}{\mathbf{S}_{\mathbf{xx}}}
      \equiv \frac{1}{n} \sum_{i=1}^{n} \mathbf{x}_i\mathbf{x}_i'
      = \frac{1}{n}\mathbf{X}'\mathbf{X}.
      \label{eq:2.3.6}
    \end{equation}
\end{enumerate}
\end{proposition}

The proof is a showcase of all the standard tricks in asymptotics.  For proving (a)
and (b), three tricks will be employed: (1)~write the object in question in terms of
sample means, (2)~apply the relevant LLN (the Ergodic Theorem in the present
context) and CLT (the ergodic stationary Martingale Differences CLT) to sample
means, and (3)~use Lemma~\ref{lem:2.4}(c) to derive the asymptotic distribution.
Proof of (c) will not be given because it is an immediate implication of ergodic
stationarity.

\begin{proof}[Proof (Parts (a) and (b))]
\begin{enumerate}
\item We first write the sampling error $\mathbf{b} - \bm{\beta}$ in terms of sample
  means.
  \begin{align}
    \mathbf{b} - \bm{\beta}
    &= (\mathbf{X}'\mathbf{X})^{-1}\mathbf{X}'\bm{\varepsilon} \notag \\
    &= \left(\frac{1}{n}\mathbf{X}'\mathbf{X}\right)^{-1}
       \left(\frac{1}{n}\mathbf{X}'\bm{\varepsilon}\right) \notag \\
    &= \left(\frac{1}{n}\sum_{i=1}^{n} \mathbf{x}_i\mathbf{x}_i'\right)^{-1}
       \left(\frac{1}{n}\sum_{i=1}^{n} \mathbf{x}_i \cdot \varepsilon_i\right)
       \notag \\
    &\equiv \mathbf{S}_{\mathbf{xx}}^{-1}\,\bar{\mathbf{g}},
    \label{eq:2.3.7}
  \end{align}
  where
  \[
    \bar{\mathbf{g}} \equiv \frac{1}{n}\sum_{i=1}^{n} \mathbf{g}_i, \quad
    \mathbf{g}_i \equiv \mathbf{x}_i \cdot \varepsilon_i.
  \]
  The sample means $\mathbf{S}_{\mathbf{xx}}$ and $\bar{\mathbf{g}}$ depend on the
  sample size $n$, although the notation does not make it explicit.

\item (Consistency)\quad Since by Assumption~\ref{ass:2.2}
  $\{\mathbf{x}_i\mathbf{x}_i'\}$ is ergodic stationary,
  $\mathbf{S}_{\mathbf{xx}} \pto \bm{\Sigma}_{\mathbf{xx}}$.  (The convergence is
  actually almost surely, but almost sure convergence implies convergence in
  probability.)  Since $\bm{\Sigma}_{\mathbf{xx}}$ is invertible by
  Assumption~\ref{ass:2.4}, $\mathbf{S}_{\mathbf{xx}}^{-1} \pto
  \bm{\Sigma}_{\mathbf{xx}}^{-1}$ by Lemma~\ref{lem:2.3}(a).  Similarly,
  $\bar{\mathbf{g}} \pto \E(\mathbf{g}_i)$ which by Assumption~\ref{ass:2.3} is
  $\mathbf{0}$.  So by Lemma~\ref{lem:2.3}(a),
  $\mathbf{S}_{\mathbf{xx}}^{-1}\bar{\mathbf{g}} \pto
  \bm{\Sigma}_{\mathbf{xx}}^{-1}\mathbf{0} = \mathbf{0}$.  Therefore,
  $\plim_{n \to \infty}(\mathbf{b} - \bm{\beta}) = \mathbf{0}$, which implies
  $\plim_{n \to \infty} \mathbf{b} = \bm{\beta}$.

\item (Asymptotic normality)\quad Rewrite \eqref{eq:2.3.7} as
  \begin{equation}
    \sqrt{n}(\mathbf{b} - \bm{\beta})
    = \mathbf{S}_{\mathbf{xx}}^{-1}(\sqrt{n}\,\bar{\mathbf{g}}).
    \label{eq:2.3.8}
  \end{equation}
  As mentioned in the statement of Assumption~\ref{ass:2.5},
  $\sqrt{n}\,\bar{\mathbf{g}} \dto N(\mathbf{0}, \mathbf{S})$.  So, by
  Lemma~\ref{lem:2.4}(c), $\sqrt{n}(\mathbf{b} - \bm{\beta})$ converges to a normal
  distribution with mean $\mathbf{0}$ and variance
  $\bm{\Sigma}_{\mathbf{xx}}^{-1}\mathbf{S}(\bm{\Sigma}_{\mathbf{xx}}^{-1})'$.
  But since $\bm{\Sigma}_{\mathbf{xx}}$ is symmetric, this expression equals
  \eqref{eq:2.3.4}. \qedhere
\end{enumerate}
\end{proof}

This result says that the distribution of $\sqrt{n}$ times the sampling error is
approximated arbitrarily well by a normal distribution when the sample size is
sufficiently large.  The natural question is how large is ``large'': how large must
the sample size be for the asymptotic approximation to be valid?  The asymptotic
result we just derived holds for all the DGPs satisfying the model assumptions.
However, the sample size needed to achieve a given measure of proximity to the
asymptotic distribution depends on the DGP.  We will partially address this issue in
the Monte Carlo experiment of this chapter.

%% ─── s^2 Is Consistent ────────────────────────────────────────────────────
\subsection*{$s^2$ Is Consistent}

We now turn to the OLS estimator, $s^2$, of the error variance.

\begin{proposition}[consistent estimation of error variance]
\label{prop:2.2}
Let $e_i \equiv y_i - \mathbf{x}_i'\mathbf{b}$ be the OLS residual for observation
$i$.  Under Assumptions~\ref{ass:2.1}--\ref{ass:2.4},
\[
  s^2 \equiv \frac{1}{n - K} \sum_{i=1}^{n} e_i^2 \pto \E(\varepsilon_i^2),
\]
provided $\E(\varepsilon_i^2)$ exists and is finite.
\end{proposition}

If we could observe the error term $\varepsilon_i$, then the obvious estimator would
be the sample mean of $\varepsilon_i^2$.  It is consistent by ergodic stationarity.
The message of Proposition~\ref{prop:2.2} is that the substitution of the OLS
residual $e_i$ for the true error term $\varepsilon_i$ does not impair consistency.
Let us go through a sketch of the proof, because knowing how to handle the
discrepancy between $\varepsilon_i$ and its estimate $e_i$ will be useful in other
contexts as well.  Since
\[
  s^2 = \frac{n}{n - K}\left(\frac{1}{n}\sum_{i=1}^{n} e_i^2\right),
\]
it suffices to prove that the sample mean of $e_i^2$,
$\frac{1}{n}\sum_i e_i^2$, converges in probability to $\E(\varepsilon_i^2)$.
The relationship between $e_i$ and $\varepsilon_i$ is given by
\begin{align}
  e_i &\equiv y_i - \mathbf{x}_i'\mathbf{b} \notag \\
  &= y_i - \mathbf{x}_i'\bm{\beta} - \mathbf{x}_i'(\mathbf{b} - \bm{\beta})
     \quad (\text{by adding and subtracting } \mathbf{x}_i'\bm{\beta}) \notag \\
  &= \varepsilon_i - \mathbf{x}_i'(\mathbf{b} - \bm{\beta}),
  \label{eq:2.3.9}
\end{align}
so that
\begin{equation}
  e_i^2 = \varepsilon_i^2
  - 2(\mathbf{b} - \bm{\beta})'\mathbf{x}_i \cdot \varepsilon_i
  + (\mathbf{b} - \bm{\beta})'\mathbf{x}_i\mathbf{x}_i'(\mathbf{b} - \bm{\beta}).
  \label{eq:2.3.10}
\end{equation}
Summing over $i$, we obtain
\begin{equation}
  \frac{1}{n}\sum_{i=1}^{n} e_i^2
  = \frac{1}{n}\sum_{i=1}^{n} \varepsilon_i^2
  - 2(\mathbf{b} - \bm{\beta})'\frac{1}{n}\sum_{i=1}^{n}
    \mathbf{x}_i \cdot \varepsilon_i
  + (\mathbf{b} - \bm{\beta})'\left(\frac{1}{n}\sum_{i=1}^{n}
    \mathbf{x}_i\mathbf{x}_i'\right)(\mathbf{b} - \bm{\beta})
  = \frac{1}{n}\sum_{i=1}^{n} \varepsilon_i^2
  - 2(\mathbf{b} - \bm{\beta})'\bar{\mathbf{g}}
  + (\mathbf{b} - \bm{\beta})'\mathbf{S}_{\mathbf{xx}}
    (\mathbf{b} - \bm{\beta}).
  \label{eq:2.3.11}
\end{equation}
The rest of the proof, which is to show that the plims of the last two terms are zero
so that $\plim \frac{1}{n}\sum_i e_i^2 = \plim \frac{1}{n}\sum_i \varepsilon_i^2$,
is left as a review question.  If you do the review question, it should be clear to
you that all that is required for the coefficient estimator is consistency; if some
consistent estimator, rather than the OLS estimator $\mathbf{b}$, is used to form
the residuals, the error variance estimator is still consistent for $\E(\varepsilon_i^2)$.
